\documentclass[11pt,a4paper]{article}

% ---------- Packages ----------
\usepackage[a4paper,margin=2.5cm]{geometry}
\usepackage[T1]{fontenc}
\usepackage[utf8]{inputenc}
\usepackage{lmodern}
\usepackage{microtype}
\usepackage{setspace}
\usepackage{titlesec}
\usepackage{enumitem}
\usepackage[numbers,sort&compress]{natbib}
\usepackage[hidelinks]{hyperref}
\usepackage{url}

% ---------- Page formatting ----------
\setstretch{1.05}
\setlength{\parindent}{0pt}
\setlength{\parskip}{4pt}

\titleformat{\section}
    {\normalfont\Large\bfseries}
    {\thesection}
    {0.8em}
    {}

\titlespacing*{\section}{0pt}{1.3em}{0.6em}

\titleformat{\subsection}
    {\normalfont\large\bfseries}
    {\thesubsection}
    {0.8em}
    {}

\titlespacing*{\subsection}{0pt}{1em}{0.4em}

\setlist[itemize]{
    leftmargin=1.6em,
    itemsep=2pt,
    topsep=3pt
}

\setlist[enumerate]{
    leftmargin=1.8em,
    itemsep=2pt,
    topsep=3pt
}

% ---------- Proposal information ----------
\newcommand{\ProposalTitle}{
    Uncertainty-Guided Agentic Reasoning: Investigating the Impact of Uncertainty Scores on LLM Agents
}

\newcommand{\StudentName}{
    Marsel Berheev m.berheev@innopolis.university
}

\newcommand{\SupervisorName}{
    Vladimir Ivanov  v.ivanov@innopolis.ru 
}

\begin{document}

% ==================================================
% Title
% ==================================================

\begin{center}
    {\LARGE\bfseries \ProposalTitle\par}

    \vspace{1.2em}

    {\large Student: \StudentName\par}

    \vspace{0.3em}

    {\large Supervisor: \SupervisorName\par}
\end{center}

\vspace{1.2em}

% ==================================================
% Abstract
% ==================================================

\begin{abstract}
Current agentic systems demonstrate outstanding performance among various types of tasks,
including code generation and validation. However due to probabilistic nature of models
and inconsistency in dateset quality, Large Language Models (LLMs), tend to hallucinate,
generating structurally correct, but logically wrong code snippets. Hallucination detections methods
vary in their methodology, however one specific subset stands out among others. Present papers demonstrated
ability of hidden states to store information about models awareness of hallucinated samples. Using this 
insights it is possible to detect, but not affect generation of hallucinations. In this paper we plan to
investigate performance of hallucination aware agentic loops. Resulting system will not only be
able to detect hallucinated samples, but also apply this knowledge for generation of correct code.
\end{abstract}

% ==================================================
% 1. Introduction
% ==================================================

\section{Introduction}



% ==================================================
% 2. Novelty and Contributions
% ==================================================

\section{Novelty and Contributions}

Main contributions of this project include:

\begin{itemize}

    \item \textbf{Hallucination Aware Agents:}
    AS a result it is planned to get an agent that automatically detect
    and fix hallucinated code snippets

    \item \textbf{Unified Uncertainty Awareness Score Interface}
    For high modularity and portability, we plan to introduce unified
    interface, that will allow to use of different implementations of UAS
    without complex adapters.

\end{itemize}

% ==================================================
% 3. State of the Art
% ==================================================

\section{State-of-the-Art}

Currently, most of papers concentrate on detection, but not
mitigation of hallucinations

\begin{itemize}
    \item \textbf{Internal Representations as Markers} - various paper proposed
    approach that allows use internal states as markers for detection of models
    uncertainty regarding generation of next token. SUch approaches as SPALMA \cite{azaria2023internalstatellmknows}, HoloScope \cite{du2024haloscopeharnessingunlabeledllm}
    and Uncertainty Quantification Heads \cite{shelmanov2025headpredictheadquestion}, have different implementations of the same 
    idea - hidden states of model store information about possible hallucination.
    \item \textbf{Middle layers are the most informative} - Many researches \cite{azaria2023internalstatellmknows}, \cite{duan2024llmsknowhallucinationempirical} demonstrated 
    that most useful information usually stored inside middle layers of models
    \item \textbf{Hallucination detection in code generation} - Some papers \cite{andriushchenko2026efficient} proposed idea
    to use such uncertainty scores for generating reliable code without hallucination 
\end{itemize}

Summing up, although we can pretty efficiently detect
hallucination, till this moment there is no proper way to use
this information in real, fully autonomous agentic applications 

% ==================================================
% 4. Goal of the Thesis
% ==================================================

\section{Goal of the Thesis}

The main goal of this research is to create hallucination aware
system, that will use information, extracted from internal states,
to detect and fix hallucinations

Specifically, the research aims to:

\begin{itemize}

    \item Define efficient and clear metrics for hallucination detection

    \item Develop an agentic loop, that will use defined metrics.

    \item Evaluate the proposed approach using appropriate
    datasets and metrics.

    \item Compare the results with existing methods
    and analyze the limitations.

\end{itemize}

% ==================================================
% 5. Detailed Work Plan and Schedule
% ==================================================

\section{Detailed Work Plan and Schedule}

The whole project planned to span two semesters. First one will be
dedicated for exploration of current metrics and wys to detect hallucinations.
Additionally it is planned to create first iteration of simple loop, using
small local LLM. The second semester will concentrate on improving results
obtained in first one

\subsection*{Semester 1}

\begin{itemize}

    \item \textbf{Literature review and problem formulation
    (September--October):}

    \begin{itemize}
        \item Review relevant scientific literature.
        \item Identify limitations of existing approaches.
        \item Define evaluation metrics and experimental setup.
    \end{itemize}

    \item \textbf{Methodology and preliminary experiments
    (November--December):}

    \begin{itemize}
        \item Select appropriate uncertainty metrics and models.
        \item Conduct preliminary experiments.
        \item Analyze the initial results.
    \end{itemize}

\end{itemize}

\subsection*{Semester 2}

\begin{itemize}

    \item \textbf{Main experiments and evaluation
    (January--March):}

    \begin{itemize}
        \item Implement the proposed methodology.
        \item Conduct experiments and compare results.
        \item Perform error analysis.
    \end{itemize}

    \item \textbf{Validation and thesis writing
    (March--April):}

    \begin{itemize}
        \item Validate the experimental findings.
        \item Discuss the results and limitations.
    \end{itemize}

    \item \textbf{Finalization and defense preparation (May):}

    \begin{itemize}
        \item Finalize the thesis.
        \item Prepare the presentation and demonstration.
    \end{itemize}

\end{itemize}

% ==================================================
% 6. Conclusion
% ==================================================

\section{Conclusion}

Summing all up, although current approach demonstrated high performance 
in detecting hallucinations, still there is no reliable and efficient way
to use this findings in real world tasks.

This paper is planned to present a solution for this problem by creating
unified interface for various hidden states based quality metrics, and
fully autonomous agentic loop with hallucination feedback

% ==================================================
% References
% ==================================================

\bibliographystyle{ieeetr}
\bibliography{refs}

\end{document}